\documentclass[10pt,a4paper]{amsart}
\usepackage[margin=19mm]{geometry}
\usepackage{amsmath,amssymb,mathtools}
\usepackage[scr=rsfs,cal=euler]{mathalfa}
\usepackage{xfrac}
\usepackage[unicode,hidelinks,bookmarks=false]{hyperref}
\newcommand{\ie}{i.e.\ }
\newcommand{\cf}{cf.\ }
\newcommand{\resp}{respectively\ }
\newcommand{\Subsection}{\subsection}
\providecommand{\nobreakdash}{\nobreak\hbox{-}}
\setlength{\emergencystretch}{2em}
\multlinegap0pt
\usepackage{csquotes}
\usepackage{amssymb}
\usepackage[shortlabels]{enumitem}
\newtheorem{lemma}{Lemma}
\newcommand{\cl}{\operatorname{cl}}
\newcommand{\rk}{\operatorname{rk}}
\title{When is the matroid Schubert variety $\mathbb{Q}$-Gorenstein?}
\author{Townsend Porcher}
\date{Source: arXiv:2608.21637v1 (CC BY 4.0)}
\begin{document}
\maketitle

\noindent\emph{Source and licence.} When is the matroid Schubert variety $\mathbb{Q}$-Gorenstein?. Version arXiv:2608.21637v1 (CC BY 4.0), distributed by arXiv under the Creative Commons Attribution 4.0 International licence, http://creativecommons.org/licenses/by/4.0/, as recorded in the arXiv licence metadata for this exact version and shown on the abstract page https://arxiv.org/abs/2608.21637v1. Full licence notice: \texttt{licenses/arxiv-2608.21637-CC-BY-4.0.txt}.\par\smallskip

\noindent\textit{Editorial scope.} Counted unit: the article's lemma lem:chow\_ring\_computations (source0.tex lines 369--387). The article's complete original proof of it is delivered with it (source0.tex lines 847--921, one \texttt{proof} environment of 75 lines, which the article places away from the statement). No mathematics is written, repaired or re-derived: the statement and the proof are the article's own lines, in the article's own notation, and no retained line is substituted or re-flowed.  Retained uncounted as the source-local context the proof needs: lines 345--349; lines 350--354; lines 355--359.  Nothing was omitted from any retained span, so the package carries no omission record.  No retained line contains a citation, so the wrapper carries no bibliography and no bibliography entry.  No retained line contains a figure environment, an image-inclusion command or drawing code, so no image is delivered and the package declares no asset dependency.  Editorial formatting only, in addition: an AMS \texttt{amsart} preamble that restates the article's own declarations and macros used by the retained lines, namely the environments \texttt{lemma} on separate counters, together with the standard packages those lines need.  The retained environments print in this document as Lemma 1 in order of appearance, whereas the article prints the same items with the numbering of its own full document.  The complete native source of the article as distributed in the arXiv e-print for this exact version is bundled unchanged as \texttt{sources/208263/source0.tex}.
\begin{equation}
\label{eq:ideal_description_1}
\tag{R1}
y_e - \sum_{F\not\ni e} x_F \quad \text{ for every element } e \in E,
\end{equation}

\begin{equation}
\label{eq:ideal_description_2}
\tag{R2}
x_{F_1}x_{F_2} \quad \text{ for every pair of incomparable proper flats } F_1, F_2, \text{ and}
\end{equation}

\begin{equation}
\label{eq:ideal_description_3}
\tag{R3}
y_ex_F \quad \text{ for every $e \in E$ and for every proper flat $F$ not containing $e$.}
\end{equation}

\begin{lemma}
\label{lem:chow_ring_computations}

Let $I = \{i_1,\dots,i_s\}$ be a (possibly empty) rank-$s$ independent set of $M_V$, and let 
\[\mathscr{F}:F_{s+1} \subsetneq F_{s+2} \subsetneq \cdots \subsetneq F_{s+t}\]
be a (possibly empty) flag of $t$ proper flats of $M_V$, such that $(I, \mathscr{F})$ is a nearly maximal compatible pair of $M_V$.
If $\mathscr{F}$ is non-empty, we have:
\begin{enumerate}
    \item[(I)] $\deg_{W_V}((x_{F_{s+1}} \cup x_{I, \mathscr{F}}) \cap [W_V]) = -1$, and
    \item[(II)] $\deg_{W_V}((x_{F_{s+k}} \cup x_{I, \mathscr{F}}) \cap [W_V]) = 0$ for all $2 \leq k \leq t$.
\end{enumerate}
Furthermore, let $a: \mathcal{L}_{\neq r} (M_V) \to \mathbb{Z}$ be a set function, and let $\alpha$ be the element
\[\sum_{F \in \mathcal{L}_{\neq r} (M_V)} a(F)x_F\]
in  $A^1(W_V)$. Then, we have:
\begin{enumerate}
    \item[(III)] $\deg_{W_V}((\alpha \cup x_{I, \mathscr{F}}) \cap [W_V]) = a(\cl(I)) - a(F_{s+1}),$
\end{enumerate}
setting $a(F_{s+1}) := 0$ if $\mathscr{F}$ is empty.
\end{lemma}

\begin{proof}[Proof of Lemma \ref{lem:chow_ring_computations}, (I)]
It follows from the definition of a nearly maximal compatible pair that $F_{s+1}$ must be a rank-$(s+1)$ flat, $F_{s+2}$ must be a rank-$(s+2)$ flat, etc. 

Choose $a,b \in E$ such that $a \in F_{s+1} \setminus \cl(I)$ and $b \in F_{s+2} \setminus F_{s+1}.$ (If $r= s+2,$ hence there is no $F_{s+2}$, then let $F_{s+2} = E.$) Due to the relations (\ref{eq:ideal_description_1}) in the Chow ring $A^*(W_V)$, we have the following equality:
\begin{equation}
\label{eq:proof_of_i_1}
y_b - y_a = \sum_{F \not\ni b}x_F - \sum_{F \not\ni a}x_F.
\end{equation}
If a proper flat $F$ contains neither $a$ nor $b$, then the element $x_F$ appears in both of the above sums. Therefore, we can rewrite Equation (\ref{eq:proof_of_i_1}) as
\begin{equation}
\label{eq:proof_of_i_2}
y_b - y_a = \sum_{F \ni a,\: F \not\ni b}x_F - \sum_{F \not\ni a,\: F \ni b}x_F.
\end{equation}
The flat $F_{s+1}$ does contain $a,$ but does not contain $b,$ so we can rewrite Equation (\ref{eq:proof_of_i_2}) as
\begin{equation*}
\label{eq:proof_of_i_3}
y_b - y_a  = \sum_{F \in S}x_F  + x_{F_{s+1}}- \sum_{F \not\ni a,\: F \ni b}x_F,
\end{equation*}
where $S$ is the collection of proper flats that contain $a$, do not contain $b,$ and are not equal to $F_{s+1}$.
Thus, we can write
\begin{equation}
\label{eq:proof_of_i_4}
x_{F_{s+1}} = y_b - y_a  - \sum_{F \in S}x_F  + \sum_{F \not\ni a,\: F \ni b}x_F.
\end{equation}
We multiply both sides of Equation (\ref{eq:proof_of_i_4}) by $x_{I, \mathscr{F}}$ and distribute to get
\begin{equation}
\label{eq:proof_of_i_5}
\begin{split}
& x_{F_{s+1}} \cdot x_{I , \mathscr{F}} = \\ & y_b \cdot x_{I , \mathscr{F}} - y_a \cdot x_{I , \mathscr{F}}  - \sum_{F \in S}x_F \cdot x_{I , \mathscr{F}}  + \sum_{F \not\ni a,\: F \ni b}x_F \cdot x_{I , \mathscr{F}}.
\end{split}
\end{equation}
We now deal with each term on the right-hand side of Equation (\ref{eq:proof_of_i_5}) one at a time. 

\noindent
\textit{First term.} The element $b\in E,$ by construction, is not a member of $F_{s+1}$. Thus, by the set of relations (\ref{eq:ideal_description_3}) in the ring $A^*(W_V)$, we have
\begin{equation*}
y_b \cdot x_{I, \mathscr{F}} = 0.
\end{equation*}

\noindent
\textit{Second term.} Let $I' = I \, \cup \, a$. Since $a \notin \cl(I)$, we have that $I'$ must be an independent set. Furthermore, since $I \subseteq F_{s+1}$ and $a \in F_{s+1}$, we have that $I' = I \cup a \subseteq F_{s+1}$. We conclude that $(I', \mathscr{F})$ is a maximal compatible pair of $M_V$. We have
\begin{equation*}
y_a \cdot x_{I, \mathscr{F}} = y_{i_1}y_{i_2}\cdots y_{i_s} y_{a} x_{F_{s+1}}x_{F_{s+2}}\cdots x_{F_{s+t}} = x_{I',\mathscr{F}}. 
\end{equation*}

\noindent
\textit{Third term.} Suppose $F$ is a member of the collection $S$. We will show that \[x_F \cdot x_{I , \mathscr{F}} = 0.\]
We will do cases based on the rank of $F$.
\begin{itemize}

\item If $\rk(F) < s$, then $F$ does not contain every element in $I$. So, by the relations (\ref{eq:ideal_description_3}), we are done.

\item If $\rk(F) = s$, then $F$ still doesn't contain every element in $I$. (If $F$ did, then $F = \cl(I)$. But $a\in F$ and $a \notin \cl(I)$.) So, by the relations (\ref{eq:ideal_description_3}), we are done.

\item If $\rk(F) = s+1$, we have that $F$ and $F_{s+1}$ are incomparable flats, as they are of the same rank but not equal. So, by the relations (\ref{eq:ideal_description_2}), we are done.

\item If $\rk(F) > s+1$, then there is a $F' \in \mathscr{F}$ with the same rank. However, they are not equal, as $b\in F_{s+2} \subseteq F'$, but $b \notin F$. So, $F$ and $F'$ are incomparable flats, and by the relations (\ref{eq:ideal_description_2}), we are done.
\end{itemize}
So we have
\begin{equation*}
\sum_{F \in S}x_F \cdot x_{I,\mathscr{F}} = 0.
\end{equation*}

\noindent
\textit{Fourth term.} Let $F$ be a proper flat that doesn't contain $a$ and contains $b$. We will show that $x_F \cdot x_{I, \mathscr{F}} = 0$. We do this by showing that $F$ and $F_{s+1}$ are incomparable flats. This can be seen by observing that $a \in F_{s+1} \setminus F$ and $b \in F \setminus F_{s+1}$. So, by the relations (\ref{eq:ideal_description_2}), we have
\begin{equation*}
\sum_{F \not\ni a,\: F \ni b}x_F \cdot x_{I,\mathscr{F}} = 0.
\end{equation*}

Putting the four terms together, we have
\begin{equation*}
x_{F_{s+1}} \cdot x_{I ,\mathscr{F}} = -x_{I' , \mathscr{F}}.
\end{equation*}
But $\deg_{W_V}(x_{I' , \mathscr{F}} \cap [W_V]) = 1,$ so we are done.
\end{proof}

\end{document}
